% ── §2 Gray-Code Encoding of Biological Alphabets ──
\section{Gray-Code Encoding of Biological Alphabets}
\label{sec:theory-gray}

\subsection{Reflected Binary Gray Codes}
\label{sec:gray-def}

A reflected binary Gray code assigns to each non-negative integer $n$ the
value $g(n) = n \oplus (n \gg 1)$, where $\oplus$ denotes bitwise exclusive-or
and $\gg$ denotes arithmetic right shift~\cite{gray1953}.  The defining
property of this encoding is that consecutive integers map to codewords that
differ in exactly one bit position, a property that follows from the identity
$g(n) \oplus g(n{+}1) = 2^{\nu_2(n+1)}$, where $\nu_2(m)$ is the two-adic
valuation of $m$.  Over a fixed width of $w$ bits, the mapping $g : \{0,
	\ldots, 2^w{-}1\} \to \{0, \ldots, 2^w{-}1\}$ is an involution (that is,
applying it twice recovers the original value), and its image forms a
Hamiltonian cycle on the $w$-dimensional hypercube $Q_w$.

The Hamming distance between two Gray-encoded values $a$ and $b$ is defined as
\begin{equation}
	d_G(a, b) = \operatorname{popcount}\!\bigl(g(a) \oplus g(b)\bigr),
	\label{eq:gray-hamming}
\end{equation}
where $\operatorname{popcount}(x)$ counts the number of set bits in the binary
representation of $x$.  This distance inherits the metric properties of the
underlying Hamming metric on $\{0,1\}^w$ and satisfies the symmetry condition
$d_G(a,b) = d_G(b,a)$ for all $a, b < 2^w$.

\subsection{DNA and RNA Encoding}
\label{sec:gray-dna}

For nucleotide sequences, we adopt the standard two-bit encoding in which
adenine (A), cytosine (C), guanine (G), and thymine/uracil (T/U) are assigned
the binary codes $00_2 = 0$, $01_2 = 1$, $11_2 = 3$, and $10_2 = 2$,
respectively.  Under this assignment, the Gray-code transformation maps A to
$0$, C to $1$, G to $2$, and T to $3$, and the transition--transversion
classification of point mutations acquires a direct interpretation in terms of
Gray-code adjacency: transitions (purine$\leftrightarrow$purine or
pyrimidine$\leftrightarrow$pyrimidine) correspond to single bit-flips, while
transversions correspond to two-bit flips.  This property was formalised in
Lean~4 as part of the verification framework described in
Section~\ref{sec:theory-lean}.

\subsection{Amino-Acid Encoding}
\label{sec:gray-amino}

For amino-acid sequences, we extend the encoding from four nucleotides to
twenty canonical residues by assigning each amino acid a five-bit code under
the He~2012 physicochemical ordering~\cite{he2012}.  The ordering groups the
twenty residues into seven classes based on hydrophobicity, charge, and
aromaticity:
\begin{center}
	\begin{tabular}{lll}
		\toprule
		\textbf{Group} & \textbf{Residues} & \textbf{He indices} \\
		\midrule
		Aliphatic      & A, I, L, V        & 0--3                \\
		Aromatic       & F, Y, W           & 4--6                \\
		Sulfur         & M, C              & 7--8                \\
		Structure      & P, G              & 9--10               \\
		Polar          & S, T, N, Q        & 11--14              \\
		Negative       & D, E              & 15--16              \\
		Positive       & H, K, R           & 17--19              \\
		\bottomrule
	\end{tabular}
\end{center}
Under this ordering, the integer code assigned to residue $a_i$ at position
$i$ is simply its He index, denoted $\mathrm{code}(a_i) \in \{0, \ldots,
	19\}$, and the corresponding five-bit Gray code is $g(\mathrm{code}(a_i))$.
The twelve unused five-bit codewords ($20 \leq x \leq 31$) are designated as
don't-care cells during Karnaugh-map minimisation, a design choice whose
formal justification is given in Section~\ref{sec:theory-kmap-qm}.

The key property that makes this encoding suitable for Karnaugh-map analysis
is that within-group pairs of residues tend to have small Gray-Hamming
distances.  Specifically, among the $\binom{20}{2} = 190$ unordered pairs of
canonical residues, 40 pairs (ordered count: 80) lie at Gray distance~1, and
these are enriched for within-group substitutions relative to what would be
expected under random assignment of codewords to residues.  This enrichment
was verified exhaustively by native computation in Lean~4 and is stated as
the theorem \texttt{cc\_he\_dist1\_ordered\_eq\_80} in the
\texttt{ContactCircuits} module.

\subsection{Gap Handling}
\label{sec:gray-gaps}

Alignment gap characters (`--') and ambiguous residue calls (`X') are mapped to
a twenty-first state with code value 20.  This state is excluded from all
statistical calculations (entropy, mutual information, coupling estimation)
but is retained in the position arrays so that column identity is preserved
across all sequences.  The retention of gaps as state~20, rather than their
deletion, constitutes the first of two critical corrections applied to our
pipeline; the deletion of gaps causes different sequences to contribute to
different raw positions at any given column index, producing spurious
variable-position calls and inflated mutual-information values.
